\chapter{Energy, Curvature, and Radiation}

\lectureoverview{The Newtonian model becomes much more revealing once two restrictions are removed. First, a comoving galaxy need not have exactly zero total energy; keeping its energy arbitrary introduces the \(a^{-2}\) term that general relativity will reinterpret as spatial curvature. Second, the cosmic source need not be pressureless matter. Photon wavelengths stretch with the scale factor, so radiation dilutes as \(a^{-4}\) and drives the law \(a(t)\propto t^{1/2}\). A universe containing both matter and radiation is therefore radiation dominated at small \(a\) and matter dominated at larger \(a\). Along the way we shall separate local Newtonian motion from global relativistic geometry and identify the questions that cannot be settled without general relativity.}

\section{What Newton Can See}

Before extending the model, we should ask why the previous Newtonian calculation worked at all. The physical universe is described by general relativity, and its spatial geometry need not be globally Euclidean. Nevertheless, every smooth curved space is approximately flat in a sufficiently small neighborhood. A small patch of the surface of the Earth looks planar; in the same way, a region of the universe much smaller than its radius of curvature can be treated as flat.

``Small'' here is relative to the curvature scale. The patch may still span billions of light-years. Within it, neighboring galaxies can be treated by Newtonian mechanics when their relative peculiar velocities are nonrelativistic and the gravitational field is weak. The Newtonian calculation is therefore not a rival to Einstein's theory. It is the local, slow-motion limit of the more complete description.

There is an immediate warning. Radiation does not move slowly. Photons, and neutrinos while they are relativistic, cannot be treated as a gas of ordinary Newtonian particles. Newtonian reasoning will continue to organize the expansion equation, but the source on its right-hand side must be enlarged from rest-mass density to energy density. This is where the approximation begins to reveal its own boundary.

\section{The Scale Factor Recalled}

Take a homogeneous distribution of galaxies and attach coordinates to a grid that moves with the average cosmic flow. A galaxy at fixed comoving coordinate \(x\) has physical distance
\begin{equation}
D(t)=a(t)x
\label{eq:c02-distance}
\end{equation}
from the chosen origin. The scale factor \(a(t)\) converts coordinate separation into physical separation.

The absolute value of \(a\) is conventional. If all comoving coordinates are rescaled by a constant, the same physical distances can be represented by the inverse rescaling of \(a\). Thus neither \(a\) nor the mass assigned to one coordinate cell is separately invariant. Ratios of scale factors and the fractional expansion rate are invariant:
\begin{equation}
\frac{a(t_2)}{a(t_1)},
\qquad
H(t)\equiv\frac{\dot a}{a}.
\label{eq:c02-hubble}
\end{equation}
Indeed, if \(a'=\alpha a\) for a fixed constant \(\alpha\), then \(\dot a'/a'=\dot a/a\).

Let \(\nu_m\) be the mass contained in one unit of comoving coordinate volume. Its numerical value changes if the grid is relabeled, but it is constant in time when ordinary matter is conserved. A comoving unit cube has physical volume \(a^3\), so the matter density is
\begin{equation}
\boxed{\rho_m(a)=\frac{\nu_m}{a^3}.}
\label{eq:c02-matter-density}
\end{equation}
This \(a^{-3}\) dilution is simply volume growth: the rest energy in the comoving box remains fixed while its volume grows.

Real galaxies are not nailed rigidly to the grid. They orbit in groups, fall into clusters, and carry small velocities relative to the average flow. The scale factor describes the smoothed motion on large scales; such local deviations are called \emph{peculiar velocities}.

\begin{classroomqa}
\audiencequestion{Must every galaxy move exactly radially away from every other galaxy?}
\lecturerresponse{Only the averaged Hubble flow has that form. A galaxy's actual motion is the sum of the common expansion and a peculiar velocity caused by nearby structure. The distinction is like a river current with eddies superposed on it: the eddies matter locally without erasing the large-scale flow. \lecturetimestamp{00:24:47--00:26:33}}
\end{classroomqa}

For pressureless matter with zero Newtonian total energy, the preceding lecture found
\begin{equation}
H^2=\frac{8\pi G}{3}\rho_m
=\frac{8\pi G\nu_m}{3a^3}.
\label{eq:c02-zero-energy-friedmann}
\end{equation}
Because the overall grid normalization is arbitrary, we may temporarily choose units for which \(8\pi G\nu_m/3=1\). On the expanding branch,
\begin{equation}
\frac{\dot a}{a}=\frac{1}{a^{3/2}},
\qquad
\dot a=\frac{1}{\sqrt a}.
\end{equation}
Inverting the derivative and integrating gives
\begin{equation}
\frac{dt}{da}=\sqrt a,
\qquad
t-t_B=\frac{2}{3}a^{3/2},
\end{equation}
where \(t_B\) is the time at which this idealized solution reaches \(a=0\). Hence
\begin{equation}
\boxed{a(t)\propto (t-t_B)^{2/3}.}
\label{eq:c02-matter-law}
\end{equation}
The universe expands forever in this critical matter-only model. Its expansion slows because \(\dot a\propto a^{-1/2}\), but \(\dot a\) does not vanish at any finite value of \(a\).

\begin{figure}[ht]
\centering
\includegraphics[width=0.68\textwidth]{lecture_02_figure_02.png}
\caption{The recap board records \(\dot a=1/\sqrt a\) and the boxed matter-dominated law \(a\propto t^{2/3}\). The intermediate algebra is partly obscured and is supplied by equations~\eqref{eq:c02-zero-energy-friedmann}--\eqref{eq:c02-matter-law}.}
\label{fig:c02-matter-recap}
\lectureframe{00:29:20.838}
\end{figure}

\begin{classroomqa}
\audiencequestion{Why was the negative square root discarded?}
\lecturerresponse{It was not discarded by the equation. Equation~\eqref{eq:c02-zero-energy-friedmann} fixes \(H^2\), not the sign of \(H\). The positive root describes expansion and the negative root describes contraction. A rock in a planet's gravitational field can be moving outward or inward while obeying the same energy equation; its initial condition selects the branch. We selected the expanding branch because that is the universe under discussion. \lecturetimestamp{00:23:09--00:24:46}}
\end{classroomqa}

The next model will still omit dark energy. That omission is deliberate: we first understand the older matter-and-radiation framework before adding the component that controls the present accelerated era. Dark matter, by contrast, is already included in the matter term because nonrelativistic dark matter also dilutes as \(a^{-3}\).

\begin{classroomqa}
\audiencequestion{Is dark energy not the dominant component today?}
\lecturerresponse{Yes. The two-component story developed here is the standard pedagogical cosmology from before the discovery of the present acceleration. Dark matter belongs with matter; dark energy is a third component to be introduced separately. Beginning with the older model lets us see what each additional term changes. \lecturetimestamp{00:28:02--00:28:45}}
\end{classroomqa}

\section{Restoring the Energy Constant}

The zero-energy choice is a special initial condition. To remove it, return to a galaxy of mass \(m\) on the boundary of a homogeneous sphere. Let its comoving radius be \(x\), so that
\begin{equation}
D=ax,
\qquad
\dot D=\dot a x.
\label{eq:c02-comoving-particle}
\end{equation}
Newton's shell theorem says that only the mass \(M(D)\) inside the sphere affects the radial acceleration; exterior spherical shells exert no net force. The galaxy's conserved Newtonian energy is
\begin{equation}
\frac{1}{2}m\dot D^{\,2}-\frac{GM(D)m}{D}=E.
\label{eq:c02-newton-energy}
\end{equation}
The factor \(m\) belongs in both terms. It is easy to omit while writing rapidly on a board, but retaining it makes clear that the final constant is the energy per unit test mass.

Homogeneity allows us to choose the representative galaxy at \(x=1\). Then \(D=a\), \(\dot D=\dot a\), and the sphere has
\begin{equation}
V_{\rm sph}=\frac{4\pi}{3}a^3,
\qquad
M(a)=\rho_m V_{\rm sph}=\frac{4\pi}{3}\rho_m a^3.
\label{eq:c02-enclosed-mass}
\end{equation}
The board uses \(V\) both for volume and, in a side annotation, for radial velocity. We shall keep the unambiguous symbols \(V_{\rm sph}\) and \(\dot D\).

\begin{figure}[ht]
\centering
\includegraphics[width=0.86\textwidth]{lecture_02_figure_03.png}
\caption{The enclosed-mass sphere and the nonzero-energy Friedmann equation. The visible labels include \(V_{\rm sph}=4\pi a^3/3\), the central mass \(M\), \(x=1\), \(D=a\), and \(\dot D=\dot a\). Part of the equation is hidden by the lecturer and is reconstructed algebraically below.}
\label{fig:c02-energy-board}
\lectureframe{00:37:59.876}
\end{figure}

\begin{figure}[ht]
\centering
\begin{tikzpicture}[scale=1.05,>=Latex]
  \draw[thick] (0,0) circle (2.6);
  \fill (0,0) circle (2.2pt) node[below left=2pt] {\(M(a)\)};
  \draw[thick,->] (0,0) -- (2.58,0);
  \fill (2.6,0) circle (2.2pt) node[right=4pt] {galaxy \(m\)};
  \node[above] at (1.30,0.04) {\(D=ax\)};
  \node[below] at (1.30,-0.04) {\(x=1\Rightarrow D=a\)};
  \draw[->,blue!60!black,thick] (2.63,0.18) -- (3.55,0.72)
    node[right] {\(\dot D=\dot a\)};
  \node at (0,-3.05) {\(M(a)=\frac{4\pi}{3}\rho_m a^3\)};
\end{tikzpicture}
\caption{A clean reconstruction of the Newtonian sphere at \(x=1\). The exterior shells cancel, so the boundary galaxy responds only to the enclosed mass.}
\label{fig:c02-sphere-redraw}
\end{figure}

Insert \(D=a\) into equation~\eqref{eq:c02-newton-energy}, divide by \(m/2\), and define the constant
\begin{equation}
C\equiv\frac{2E}{m}.
\end{equation}
We obtain
\begin{equation}
\dot a^{\,2}-\frac{2GM(a)}{a}=C.
\label{eq:c02-energy-a}
\end{equation}
Dividing by \(a^2\) and using equation~\eqref{eq:c02-enclosed-mass} gives
\begin{equation}
\left(\frac{\dot a}{a}\right)^2-
\frac{8\pi G}{3}\rho_m=\frac{C}{a^2},
\label{eq:c02-energy-form}
\end{equation}
or, equivalently,
\begin{equation}
\boxed{H^2=\frac{8\pi G}{3}\rho_m+\frac{C}{a^2}.}
\label{eq:c02-friedmann-C}
\end{equation}
This is the first Friedmann equation for pressureless matter, written in the Newtonian energy convention. The new \(a^{-2}\) term remembers whether the representative galaxy was launched above, at, or below escape energy.

Substituting \(\rho_m=\nu_m/a^3\), it is convenient to define
\begin{equation}
A_m\equiv\frac{8\pi G\nu_m}{3}>0,
\end{equation}
so that
\begin{equation}
H^2=\frac{A_m}{a^3}+\frac{C}{a^2},
\qquad
\dot a^{\,2}=\frac{A_m}{a}+C.
\label{eq:c02-energy-branches}
\end{equation}
The equation is more informative when read by limits than when immediately integrated.

\subsection{Positive Energy}

If \(C>0\), the matter term \(A_m/a^3\) dominates at small \(a\). The early expansion therefore retains the matter law \(a\propto t^{2/3}\). At large \(a\), however, the \(C/a^2\) term dominates in \(H^2\), or equivalently
\begin{equation}
\dot a^{\,2}\longrightarrow C.
\end{equation}
The expansion approaches constant speed,
\begin{equation}
a(t)\sim\sqrt C\,t.
\label{eq:c02-positive-late}
\end{equation}
This is the cosmological counterpart of throwing an apple faster than escape velocity. Gravity continually reduces its speed, but far away the apple retains a nonzero asymptotic velocity.

\subsection{Zero Energy}

For \(C=0\), equation~\eqref{eq:c02-energy-branches} returns the critical solution \(a\propto t^{2/3}\). It expands without bound while \(\dot a\) tends to zero. This is the exact escape-energy case.

\subsection{Negative Energy}

If \(C=-|C|<0\), then
\begin{equation}
H^2=\frac{A_m}{a^3}-\frac{|C|}{a^2}
=\frac{A_m-|C|a}{a^3}.
\label{eq:c02-negative-energy}
\end{equation}
The scale factor cannot exceed
\begin{equation}
a_{\max}=\frac{A_m}{|C|},
\label{eq:c02-amax}
\end{equation}
because the right-hand side reaches zero there. The expanding branch arrives at \(a_{\max}\) with \(\dot a=0\); the continuation is the contracting branch. In the Newtonian picture this is an apple launched below escape velocity: it rises, turns, and falls.

\begin{classroomqa}
\audiencequestion{How can the model recollapse when the right-hand side equals \(H^2\) and therefore cannot be negative?}
\lecturerresponse{It never becomes negative. During expansion the right-hand side decreases to zero at \(a_{\max}\). There \(H=0\). The solution then changes from the positive square root to the negative square root and retraces the allowed region with \(\dot a<0\). A tempting but incorrect continuation would let the expression pass through zero. It cannot: zero is the boundary, not a value through which a real \(H\) can continue. \lecturetimestamp{00:48:43--00:53:03}}
\end{classroomqa}

Thus the three Newtonian energies give three qualitative histories: recollapse for \(C<0\), critical escape for \(C=0\), and eternal expansion with asymptotically constant \(\dot a\) for \(C>0\). The surprising next step is that general relativity assigns these same cases a geometric meaning.

\section{From Energy to Curvature}

Move the energy term in equation~\eqref{eq:c02-friedmann-C} to the geometric side and define
\begin{equation}
k\equiv-C.
\label{eq:c02-k-definition}
\end{equation}
Then
\begin{equation}
\boxed{H^2+\frac{k}{a^2}=\frac{8\pi G}{3}\rho.}
\label{eq:c02-friedmann-k}
\end{equation}
The board alternates among \(C\), \(K\), and \(\kappa\), and a sign is absorbed when moving the term across the equality. Equation~\eqref{eq:c02-k-definition} fixes our convention: positive Newtonian energy corresponds to \(k<0\), zero energy to \(k=0\), and negative Newtonian energy to \(k>0\).

\begin{figure}[ht]
\centering
\includegraphics[width=0.76\textwidth]{lecture_02_figure_04.png}
\caption{The Friedmann equation arranged as geometry on the left and source on the right. The frame supports this conceptual division; it does not display the full tensor Einstein equation.}
\label{fig:c02-geometry-source}
\lectureframe{00:55:22.810}
\end{figure}

This arrangement is a miniature version of Einstein's equation. The left side describes geometry: expansion through \(H\), and spatial curvature through \(k/a^2\). The right side describes the material source. In Newtonian mechanics that source was mass density. In relativity it is energy and momentum, represented fully by the stress-energy tensor. For a homogeneous and isotropic universe the first Friedmann equation depends on the total energy density \(\rho\).

If ordinary units are retained, a mass density contributes through its rest energy \(\rho_{\rm mass}c^2\). From here onward we set \(c=1\), so mass and energy use the same units. The symbol \(\rho\) may then include rest mass, kinetic energy, radiation, and any other homogeneous energy component.

\begin{figure}[ht]
\centering
\includegraphics[width=0.88\textwidth]{lecture_02_figure_05.png}
\caption{Side-by-side board forms of the Friedmann equation: a curvature form using \(\rho\), and a matter form obtained after \(\rho_m=\nu_m/a^3\). The notation changes on the board; equation~\eqref{eq:c02-k-definition} keeps the sign convention fixed here.}
\label{fig:c02-curvature-comparison}
\lectureframe{00:57:52.659}
\end{figure}

The geometric interpretation is global. In the relativistic model, \(k=0\) describes flat spatial slices, \(k>0\) positively curved spherical slices, and \(k<0\) negatively curved hyperbolic slices. The Newtonian derivation reproduced the same evolution equation without possessing these global geometries. Its nonzero energy constant is therefore a remarkably accurate mimic of curvature, but not a derivation of curved space itself.

\begin{classroomqa}
\audiencequestion{How can Newtonian mechanics and general relativity produce the same expansion equation if their geometries are different?}
\lecturerresponse{General relativity includes the flat case, and locally every smooth geometry is approximately flat. More remarkably, the Newtonian energy constant enters the evolution equation in exactly the place occupied by relativistic spatial curvature. The dynamical formula agrees, while the global interpretation does not: Euclidean Newtonian space has not secretly become a sphere or a hyperbolic space. The three geometries will have to be constructed within general relativity. \lecturetimestamp{01:20:59--01:24:20}}
\end{classroomqa}

\section{Radiation in an Expanding Box}

We now change the contents of the universe. Consider one comoving unit cube. Its physical volume is \(a^3\). For nonrelativistic matter, the number of particles in that box is fixed and each particle's rest energy is approximately fixed, giving \(\rho_m\propto a^{-3}\). For radiation there is another effect: each photon loses energy as the universe expands.

A photon of wavelength \(\lambda\) has energy
\begin{equation}
E_\gamma=\frac{hc}{\lambda},
\qquad
E_\gamma\propto\lambda^{-1}.
\label{eq:c02-photon-energy}
\end{equation}
Under homogeneous expansion, a freely propagating wavelength stretches with the scale factor:
\begin{equation}
\frac{\lambda(t_2)}{\lambda(t_1)}
=\frac{a(t_2)}{a(t_1)},
\qquad
\lambda\propto a.
\label{eq:c02-redshift}
\end{equation}
This is cosmological redshift. It implies
\begin{equation}
E_\gamma\propto a^{-1}.
\label{eq:c02-photon-redshift}
\end{equation}
If photon number in the comoving box is conserved, its total radiation energy falls as \(a^{-1}\). Dividing by its physical volume \(a^3\) gives
\begin{equation}
\boxed{\rho_r(a)=\frac{\nu_r}{a^4}.}
\label{eq:c02-radiation-density}
\end{equation}
Here \(\nu_r\) is the constant amplitude of the freely streaming radiation density. Radiation has one more inverse power of \(a\) than matter: three powers from spatial dilution and one from the redshift of every photon.

The guitar-string analogy gives useful intuition. A mode fitted to a string has an integer number of half-wavelengths. If the string is lengthened slowly, that integer cannot change continuously; instead the wavelength adjusts to the new length. Likewise, for a wave on a slowly expanding ring, the number of nodes is an adiabatic invariant, so its wavelength follows the circumference. A general wave can be decomposed into such modes.

This is not a substitute for solving Maxwell's equations in an expanding spacetime. It is the controlled intuition behind equation~\eqref{eq:c02-redshift}: slow geometric expansion changes the physical wavelength while preserving the mode label.

\begin{classroomqa}
\audiencequestion{Would this argument require a conducting or reflecting box? What happens in empty expanding space with no physical walls?}
\lecturerresponse{The walls are a bookkeeping device, not a claim that the universe is enclosed by a conductor. A full proof follows from wave propagation in the expanding metric. For the present argument we use the observed and relativistically derived fact that freely propagating wavelengths scale with \(a\). The box helps count photons and physical volume; it does not cause the cosmological redshift. \lecturetimestamp{01:02:21--01:04:37}}
\end{classroomqa}

\begin{classroomqa}
\audiencequestion{Why should a wavelength track the grid rather than keep a fixed physical length?}
\lecturerresponse{Consider a mode on a ring whose circumference changes adiabatically. Its integer node count cannot drift continuously, so the allowed wavelength must change with the ring. Expanding a general wave into Fourier modes gives the same conclusion mode by mode. This supplies the needed adiabatic intuition; the exact statement belongs to field theory in curved spacetime. \lecturetimestamp{01:34:18--01:38:33}}
\end{classroomqa}

\subsection{The Radiation-Dominated Solution}

With curvature omitted for the moment, insert equation~\eqref{eq:c02-radiation-density} into Friedmann's equation:
\begin{equation}
H^2=\frac{8\pi G\nu_r}{3a^4}
\equiv\frac{A_r}{a^4},
\qquad A_r>0.
\label{eq:c02-radiation-friedmann}
\end{equation}
On the expanding branch,
\begin{equation}
\frac{\dot a}{a}=\frac{\sqrt{A_r}}{a^2},
\qquad
\dot a=\frac{\sqrt{A_r}}{a}.
\end{equation}
Therefore
\begin{equation}
a\,da=\sqrt{A_r}\,dt,
\qquad
\frac{a^2}{2}=\sqrt{A_r}(t-t_B),
\end{equation}
and
\begin{equation}
\boxed{a(t)\propto(t-t_B)^{1/2}.}
\label{eq:c02-radiation-law}
\end{equation}
The distinction between the exponents \(1/2\) and \(2/3\) is visually modest but physically decisive. Expansion rate, temperature, reaction freeze-out, and horizon size in the early universe all depend on which component dominates.

\section{A Universe of Matter and Radiation}

For separately conserved matter and radiation, their densities add:
\begin{equation}
\rho(a)=\rho_m(a)+\rho_r(a).
\end{equation}
Ignoring curvature for this comparison,
\begin{equation}
\boxed{H^2=\frac{C_m}{a^3}+\frac{C_r}{a^4},}
\label{eq:c02-mixed-friedmann}
\end{equation}
where \(C_m,C_r>0\). The board first labels the two coefficients \(C_1\) and \(C\); the second symbol is corrected verbally to a distinct constant. The notation in equation~\eqref{eq:c02-mixed-friedmann} makes that correction explicit.

\begin{figure}[ht]
\centering
\includegraphics[width=0.74\textwidth]{lecture_02_figure_06.png}
\caption{The mixed matter-and-radiation equation on the board. The written second coefficient is subsequently corrected to a radiation-specific constant, denoted \(C_r\) in equation~\eqref{eq:c02-mixed-friedmann}.}
\label{fig:c02-mixed-board}
\lectureframe{01:11:37.997}
\end{figure}

At small \(a\), the radiation term wins because \(a^{-4}\) grows faster than \(a^{-3}\):
\begin{equation}
a\ll a_{\rm eq}
\quad\Longrightarrow\quad
H^2\simeq\frac{C_r}{a^4},
\qquad
a(t)\propto t^{1/2}.
\end{equation}
At larger \(a\), matter wins:
\begin{equation}
a\gg a_{\rm eq}
\quad\Longrightarrow\quad
H^2\simeq\frac{C_m}{a^3},
\qquad
a(t)\propto t^{2/3}.
\end{equation}
The equality scale follows directly by setting the two terms equal:
\begin{equation}
a_{\rm eq}=\frac{C_r}{C_m}.
\label{eq:c02-equality-scale}
\end{equation}
This equality is not the same event as photon decoupling. Equality asks when the two energy densities match; decoupling asks when interactions cease to keep matter and radiation in thermal contact.

The radiation sector includes photons and any species moving ultrarelativistically. Very light neutrinos behave this way at sufficiently high temperature, and a background of gravitons would scale in the same manner. Once a massive species becomes nonrelativistic, its dominant rest-mass contribution begins to scale like matter. Cold dark matter belongs in \(C_m/a^3\), not in the radiation term.

\begin{classroomqa}
\classquestion{Which term in equation~\eqref{eq:c02-mixed-friedmann} dominates first, and which dominates later?}
\lecturerresponse{For small \(a\), \(a^{-4}\) outruns \(a^{-3}\), so radiation controls the early expansion and gives \(a\propto t^{1/2}\). As \(a\) grows, the radiation-to-matter ratio falls as \(1/a\); matter then controls the expansion and gives \(a\propto t^{2/3}\). \lecturetimestamp{01:11:46--01:13:38}}
\end{classroomqa}

\begin{classroomqa}
\audiencequestion{Are the relative amounts of matter and radiation being assumed constant?}
\lecturerresponse{Their densities are not constant, but the coefficients \(C_m\) and \(C_r\) can be treated as constant after the components cease exchanging energy efficiently and particle creation or annihilation is negligible. Matter then carries a conserved comoving rest mass and free radiation a conserved comoving photon number, with redshift supplying its extra factor of \(a^{-1}\). During epochs of strong interaction or species annihilation, the bookkeeping must be refined. \lecturetimestamp{01:13:39--01:14:53}}
\end{classroomqa}

\begin{classroomqa}
\audiencequestion{When did matter and radiation stop interacting strongly?}
\lecturerresponse{A rough order-of-magnitude answer is \(10^5\) years for photons to decouple from ordinary matter. The precise recombination history is a later calculation, and it should not be confused with matter--radiation equality. The point needed here is only that, after decoupling, the two components may be followed approximately separately. \lecturetimestamp{01:27:29--01:27:46}}
\end{classroomqa}

\begin{classroomqa}
\audiencequestion{How can we know that an early \(a\propto t^{1/2}\) era occurred before there were galaxies to trace the expansion?}
\lecturerresponse{We infer the expansion history from the physical records it leaves. Primordial nucleosynthesis depends on the early competition between reaction rates and \(H(t)\); the microwave background records the later thermal state and perturbations. Galaxy observations make the matter-era behavior more direct, while the radiation era is reconstructed from this early-universe evidence. \lecturetimestamp{01:27:46--01:30:10}}
\end{classroomqa}

The matter-plus-radiation model is deliberately incomplete. It is the classic framework in which the universe passes from radiation domination to matter domination. Present accelerated expansion requires another source, dark energy, which will add a qualitatively different term.

\section{Where the Redshifted Energy Goes}

The decrease \(E_\gamma\propto a^{-1}\) raises an unavoidable question: where has the photon energy gone? In a literal expanding box with reflecting walls, radiation exerts pressure and performs work on the walls. The thermodynamic statement is
\begin{equation}
dE=-p\,dV.
\label{eq:c02-pdv}
\end{equation}
As \(V\) increases, the radiation energy decreases. The lost energy may appear as kinetic energy, strain, or heat in whatever moves the walls.

The comoving box has no material walls. It is an imaginary boundary moving with the average expansion. Photons cross it, but in a homogeneous radiation bath the statistical flux entering balances the flux leaving. The box remains a valid bookkeeping volume even though it is not a container. The pressure-work law still describes how the energy in a comoving volume evolves.

For radiation \(p=\rho/3\), and \(E=\rho V\). Equation~\eqref{eq:c02-pdv} then gives
\begin{align}
d(\rho V)&=-\frac{\rho}{3}\,dV,\\
V\,d\rho&=-\frac{4}{3}\rho\,dV,\\
\frac{d\rho}{\rho}&=-\frac{4}{3}\frac{dV}{V}.
\end{align}
Since \(V\propto a^3\), integration yields \(\rho\propto a^{-4}\), agreeing with the wavelength argument. This thermodynamic derivation also makes clear that the energy of each photon, not merely the number density, decreases.

\begin{classroomqa}
\audiencequestion{Where does the energy of the redshifted photons go? Is total cosmic energy conserved?}
\lecturerresponse{With physical walls the radiation does work on them, so the transfer is ordinary mechanical bookkeeping. For the universe as a whole the question is more subtle. Global energy conservation normally follows from time-translation symmetry, but an expanding spacetime is not generally invariant under time translation and need not possess a unique conserved total energy. One can build useful Newtonian analogies involving the energy of expansion, but a definitive global statement requires general relativity. Both descriptions are useful here, but no final global verdict follows from the Newtonian model. \lecturetimestamp{01:16:51--01:20:58}}
\end{classroomqa}

\begin{classroomqa}
\audiencequestion{If the comoving box is imaginary and photons cross its boundary, how can pressure work still be a useful picture?}
\lecturerresponse{Homogeneity supplies the missing bookkeeping. On average, every photon leaving one comoving volume is replaced statistically by one entering from its neighbor. The local fluxes cancel in the mean, while the physical volume and wavelengths change. A reflecting box makes the work mechanism concrete; the comoving box carries the same averaged conservation law without literal walls. \lecturetimestamp{01:24:22--01:27:28}}
\end{classroomqa}

\begin{classroomqa}
\audiencequestion{Does expansion merely reduce photon density, or does it reduce the energy of every photon?}
\lecturerresponse{Both occur. The number per physical volume falls as \(a^{-3}\), and cosmological redshift lowers each photon energy by another factor \(a^{-1}\). This is the radiative analogue of adiabatic cooling. An ordinary gas also cools while doing work during slow expansion, although a sudden free expansion is a different thermodynamic process. \lecturetimestamp{01:42:30--01:45:55}}
\end{classroomqa}

\section{Recession, Horizons, and the Limit of the Model}

Hubble's law assigns the recession rate
\begin{equation}
v_{\rm rec}=HD
\label{eq:c02-recession}
\end{equation}
to an object at proper distance \(D\) on a chosen cosmic-time slice. For sufficiently large \(D\), equation~\eqref{eq:c02-recession} can exceed \(c\). This does not violate special relativity, because \(v_{\rm rec}\) is not the velocity of a nearby object measured in one local inertial frame. It compares widely separated comoving worldlines in a changing geometry. Every local observer still measures light at speed \(c\), and no local material object passes another faster than light.

\begin{classroomqa}
\audiencequestion{Can a sufficiently distant galaxy recede faster than light?}
\lecturerresponse{Yes, its cosmological recession rate \(HD\) can exceed \(c\). No local object then overtakes a neighboring observer faster than light, and the result cannot be used for superluminal signaling. It is a statement about the changing proper distance between distant comoving worldlines, not an ordinary special-relativistic velocity measured at one event. \lecturetimestamp{01:30:17--01:32:13}}
\end{classroomqa}

Superluminal recession by itself does not decide whether light emitted now will ever reach us. That is a question about the entire future expansion history. In a decelerating matter-dominated model, regions currently receding faster than light can later become observable because the relevant horizon evolves. If accelerated expansion persists, an event horizon may permanently hide sufficiently distant events.

\begin{classroomqa}
\audiencequestion{Do galaxies disappear once their recession rate exceeds the speed of light?}
\lecturerresponse{Not simply. A Hubble-radius crossing and a true event horizon are different notions. In a decelerating model, light from a currently superluminally receding region may still reach us later. Persistent dark-energy domination changes the global causal structure and can create an ultimate horizon. Deciding which signals arrive requires the spacetime geometry, not a snapshot of \(HD\). \lecturetimestamp{01:32:14--01:34:18}}
\end{classroomqa}

Special-relativistic language can also mislead when applied globally. A particle's invariant rest mass does not increase with speed. Its energy and momentum depend on the observer's local frame. Massive particles cannot be accelerated through the local speed of light, but the recession of a distant comoving galaxy is not its passage through a single local frame.

\begin{classroomqa}
\audiencequestion{Does a rapidly receding galaxy acquire a larger relativistic mass?}
\lecturerresponse{It is cleaner not to use ``relativistic mass.'' The galaxy has an invariant rest mass, while its energy and momentum are frame dependent. Moreover, a recession rate assigned across cosmological distance is not the local velocity appearing in the special-relativistic formula. These concepts agree only after the distant events are compared using the spacetime metric. \lecturetimestamp{01:38:34--01:40:05}}
\end{classroomqa}

\begin{classroomqa}
\audiencequestion{What happens to the time and space axes when a distant recession rate is superluminal?}
\lecturerresponse{There is no single global Lorentz frame whose tilted axes settle the question. Special relativity compares inertial frames in flat spacetime; cosmology requires a metric that relates local frames over an extended curved spacetime. Newtonian reasoning remains reliable in a small patch, but the global construction must wait for general relativity. \lecturetimestamp{01:40:07--01:42:30}}
\end{classroomqa}

This is the natural stopping point. Newtonian energy conservation has produced the correct scale-factor equation for matter and has even anticipated the curvature term. Wave mechanics and thermodynamics have supplied the radiation scaling. But spatial geometry, global energy, and horizons are not local Newtonian questions. They require the metric description that follows.

\section{Summary}

The matter-only critical solution obeys
\begin{equation}
\rho_m\propto a^{-3},
\qquad
a(t)\propto t^{2/3}.
\end{equation}
Restoring the Newtonian energy per unit mass gives
\begin{equation}
H^2=\frac{8\pi G}{3}\rho_m+\frac{C}{a^2}.
\end{equation}
Positive \(C\) yields eternal expansion with \(a\sim t\) at late times, \(C=0\) gives critical escape, and negative \(C\) yields a maximum scale factor followed by recollapse. Defining \(k=-C\) places the same term on the geometric side of the relativistic Friedmann equation.

Radiation differs from matter because expansion changes both number density and energy per quantum:
\begin{equation}
\begin{aligned}
\lambda&\propto a,
&\qquad E_\gamma&\propto a^{-1},\\
\rho_r&\propto a^{-4},
& a(t)&\propto t^{1/2}.
\end{aligned}
\end{equation}
The two-component equation
\begin{equation}
H^2=\frac{C_m}{a^3}+\frac{C_r}{a^4}
\end{equation}
therefore describes a radiation-dominated early universe followed by a matter-dominated era. Dark energy, spatial geometry, and the causal structure of horizons remain outside this model, exactly where the Newtonian approximation tells us to hand the problem over to general relativity.
